% TEMPLATE-MILC-v3.tex - plantilla maestra UNICA de las guias MILC v3.
%
% La usan las tres familias de guias, cada una con su driver:
%   · Grados 6-11  -> scripts/build-guias-g11.py
%   · Bebras       -> scripts/build-guias-bebras.py
%   · Semillero    -> scripts/build-guias-semillero.py
%
% Antes habia tres copias casi identicas (~400 lineas iguales, 6 distintas):
% cada mejora habia que aplicarla tres veces, y en la practica no se hacia
% (el motor de recursos vivia solo en dos de ellas). Lo que varia entre
% familias esta parametrizado en 6 placeholders que cada driver llena
% (se nombran aqui SIN su sintaxis real: el builder reemplaza por texto plano,
% tambien dentro de los comentarios, y un valor multilinea rompe el preambulo):
%
%   HEADER_IZQ        encabezado izquierdo de pagina
%   HEADER_DER        encabezado derecho de pagina
%   PORTADA_ETIQUETA  rotulo sobre el numero grande (GRADO/MOMENTO/MODULO)
%   PORTADA_SERIE     linea de serie de la portada
%   PORTADA_ID        identificador de la guia en la portada
%   PORTADA_CIRCULO   el circulito de grado (°), o vacio si no aplica
%
% El resto de claves las documenta content/guias/_SCHEMA.md. El script valida
% que no queden placeholders sin reemplazar antes de escribir el archivo final.

\documentclass[11pt,letterpaper]{article}
\usepackage[margin=2.0cm]{geometry}
\usepackage[table]{xcolor}
\usepackage{fontspec}
\usepackage[spanish]{babel}
\usepackage{tikz}
% Librerías TikZ para los diagramas de las guías (bloque `recursos:`):
%  · babel  → babel en español vuelve activos caracteres como «>», y TikZ los usa
%             en sus opciones (>=stealth, ->). Sin esta librería, cualquier
%             diagrama con flechas revienta con «\language@active@arg> has an
%             extra }». Debe ir ANTES de que se dibuje nada.
%  · positioning        → sintaxis «below left=2pt and 3pt».
%  · decorations.pathreplacing → llaves ({brace}) para señalar rangos.
\usetikzlibrary{babel,positioning,decorations.pathreplacing}
\usepackage{graphicx}
% Circuitos: el trabajo de electrónica se hace hoy con micro:bit y sus sensores
% integrados (MakeCode, sin cableado), así que no se necesitan esquemáticos.
% Si algún día entran montajes con protoboard, basta descomentar esta línea:
% los diagramas .tex/.tikz declarados en `recursos:` ya se incrustan solos.
% \usepackage{circuitikz}
\usepackage[most]{tcolorbox}
\usepackage{tabularx}
\usepackage{array}
\usepackage{fancyhdr}
\usepackage{titlesec}
\usepackage[document]{ragged2e}  % bandera a la derecha en todo el cuerpo
\usepackage{enumitem}
\usepackage{microtype}

% Helvetica Neue no trae la flecha U+2192, asi que cada «→» escrito literalmente
% en el YAML se compilaba a NADA, en silencio (462 flechas en 61 guias: rutas del
% tipo «paleta → tipografia → lockup» salian rotas). Mapeamos el caracter a su
% version matematica, que si tiene glifo. Asi el autor puede seguir escribiendo
% «→» comodamente en el YAML.
\usepackage{newunicodechar}
\newunicodechar{→}{\ensuremath{\rightarrow}}
% Mismo problema con los simbolos de marca y vinetas: el «✓» de «una palomita ✓»
% salia como cuadrito vacio en el PDF de todas las guias que lo usaban.
\usepackage{pifont}
\newunicodechar{✓}{\ding{51}}
\newunicodechar{✗}{\ding{55}}
\newunicodechar{✔}{\ding{52}}
\newunicodechar{•}{\textbullet}
\newunicodechar{≥}{\ensuremath{\geq}}
\newunicodechar{≤}{\ensuremath{\leq}}

\setmainfont{Helvetica Neue}
\setsansfont{Helvetica Neue}
\setlength{\parindent}{0pt}
\setlength{\parskip}{6pt}
% Sin particion de palabras: en 6.º-7.º una palabra cortada por guion al final
% de linea («Comuni-cacion») frena la lectura mucho mas que una linea corta.
\hyphenpenalty=10000
\exhyphenpenalty=10000
\pretolerance=2000
\tolerance=3000
\setlength{\emergencystretch}{3em}
\linespread{1.10}
\renewcommand{\arraystretch}{1.25}
\setlist{leftmargin=5mm,itemsep=1mm,topsep=1mm}
\newcolumntype{Y}{>{\RaggedRight\arraybackslash}X}

\definecolor{milcMagenta}{HTML}{E6007E}
\definecolor{milcVino}{HTML}{5A0038}
\definecolor{milcTurquesa}{HTML}{0093A5}
\definecolor{milcVerde}{HTML}{58A923}
\definecolor{milcMostaza}{HTML}{E5B400}
\definecolor{milcOcre}{HTML}{B86600}
\definecolor{milcNegro}{HTML}{241816}
\definecolor{milcGris}{HTML}{F7F7F4}
\definecolor{milcLinea}{HTML}{E8E2DD}
\definecolor{milcGradePrimary}{HTML}{0066FF}
\definecolor{milcGradeDark}{HTML}{003D99}
\definecolor{milcGradeSoft}{HTML}{D6E8FF}
\definecolor{milcGradeAccent}{HTML}{CFE7FF}
\definecolor{milcGradeText}{HTML}{FFFFFF}
\color{milcNegro}

\pagestyle{fancy}
\fancyhf{}
\lhead{\small Tecnología e Informática · Grado 6}
\rhead{\small Guía 8 · MILC v3}
\cfoot{\small\thepage}
\renewcommand{\headrulewidth}{0pt}

\newtcolorbox{titlebox}[1]{
  enhanced,colback=#1,colframe=#1,arc=7mm,boxrule=0pt,
  left=7mm,right=7mm,top=3mm,bottom=3mm,
  before skip=8pt,after skip=8pt,
  fontupper=\sffamily\bfseries\Large\color{white}
}

\newtcolorbox{softbox}[3]{
  enhanced,colback=#2,colframe=#1,borderline west={4pt}{0pt}{#1},
  arc=2mm,boxrule=.4pt,left=4mm,right=4mm,top=3mm,bottom=3mm,
  before skip=6pt,after skip=8pt,title={#3},coltitle=white,
  colbacktitle=#1,fonttitle=\sffamily\bfseries,fontupper=\RaggedRight,
  attach boxed title to top left={xshift=3mm,yshift=-2mm},
  boxed title style={arc=2mm, boxrule=0pt}
}

\newtcolorbox{drawbox}[2]{
  enhanced,colback=white,colframe=#1,arc=4mm,boxrule=1pt,
  left=5mm,right=5mm,top=4mm,bottom=4mm,before skip=8pt,after skip=8pt,
  title={#2},coltitle=white,colbacktitle=#1,fonttitle=\sffamily\bfseries,
  fontupper=\RaggedRight,
  attach boxed title to top left={xshift=4mm,yshift=-2mm},
  boxed title style={arc=3mm, boxrule=0pt}
}

\newtcolorbox{infoband}[1]{
  enhanced,colback=#1!8,colframe=#1!50,arc=2mm,boxrule=.5pt,
  left=4mm,right=4mm,top=2mm,bottom=2mm,before skip=4pt,after skip=4pt,
  fontupper=\small\RaggedRight
}

% Puente narrativo entre secciones (contrato editorial Conéctate).
% Da continuidad: "Acabas de X. Ahora vas a Y porque Z."
% Visualmente: banda izquierda en color del grado, texto en itálica pequeña.
\newtcolorbox{puentebox}{
  enhanced,colback=milcGradeSoft,colframe=milcGradeDark,
  borderline west={3pt}{0pt}{milcGradeDark},
  arc=0mm,boxrule=0pt,left=5mm,right=4mm,top=2.5mm,bottom=2.5mm,
  before skip=4pt,after skip=8pt,
  fontupper=\itshape\small\color{milcNegro}\RaggedRight
}
% La flecha va en modo matematico a proposito: Helvetica Neue NO tiene el glifo
% U+2192, asi que \textrightarrow se compilaba a NADA («Missing character: There
% is no → in font Helvetica Neue») y el puente salia sin flecha. En modo
% matematico la flecha viene de la fuente matematica y si se dibuja.
\newcommand{\puente}[1]{%
  \begin{puentebox}%
    \textcolor{milcGradeDark}{$\rightarrow$}\hspace{2mm}#1%
  \end{puentebox}%
}

\newcommand{\linead}[1]{\noindent\color{milcLinea}\rule{#1}{0.5pt}\color{milcNegro}}
\newcommand{\respuesta}[1]{\par\vspace{2mm}\foreach \n in {1,...,#1}{\noindent\color{milcLinea}\rule{\linewidth}{0.5pt}\par\vspace{5mm}}\color{milcNegro}}
\newcommand{\checkbox}{\textcolor{milcGradeDark}{\fbox{\phantom{x}}}\hspace{5pt}}

\newcommand{\phasecard}[4]{
\begin{tcolorbox}[enhanced,colback=#3,colframe=#2,arc=3mm,boxrule=.5pt,height=3.05cm,valign=center,halign=center,left=1.5mm,right=1.5mm,top=2mm,bottom=2mm]
{\sffamily\bfseries\fontsize{22}{24}\selectfont\textcolor{#2}{#1}}\par
{\sffamily\bfseries\small #4}
\end{tcolorbox}
}

% ── Piezas del rediseno 2026 (legibilidad 6.º-11.º) ───────────────────
% Banda de estacion: reemplaza el titlebox macizo. Titulo a la izquierda,
% dato practico (tiempo/modalidad) a la derecha, en una sola linea.
\newcommand{\milcestacion}[2]{%
\begin{tcolorbox}[enhanced,colback=#1,colframe=#1,arc=2mm,boxrule=0pt,
  left=5mm,right=5mm,top=2.6mm,bottom=2.6mm,before skip=11pt,after skip=7pt]
{\sffamily\bfseries\fontsize{15}{18}\selectfont\color{white}#2}
\end{tcolorbox}}

% Tarjeta de paso numerada: sustituye las tablas «Pilar / Aplicacion», que
% eran formato de consulta docente, no de estudiante. El numero va en un
% circulo sobre el borde para que la secuencia se lea de un vistazo.
\newcommand{\milcpaso}[3]{%
\begin{tcolorbox}[enhanced,colback=white,colframe=milcLinea,arc=1.5mm,boxrule=.9pt,
  left=14mm,right=4mm,top=3mm,bottom=3mm,before skip=4pt,after skip=4pt,
  fontupper=\RaggedRight,
  overlay={\node[anchor=north west,circle,fill=#2,text=white,inner sep=0pt,
    minimum size=7.4mm,font=\sffamily\bfseries\small]
    at ([xshift=3.6mm,yshift=-3.2mm]frame.north west) {#1};}]
#3
\end{tcolorbox}}

% Casilla marcable de tamano util con lapiz (la anterior era \fbox{\phantom{x}},
% de alto variable segun el texto que la seguia).
\newcommand{\milcmarca}{\framebox[3.8mm]{\rule{0pt}{3.4mm}}}

% Fila marcable de lista suelta (checks de la estacion 1).
\newcommand{\milccheckrow}[1]{%
\par\vspace{1.8mm}\noindent
\makebox[6.4mm][l]{\raisebox{-.55mm}{\textcolor{milcNegro}{\milcmarca}}}%
\parbox[t]{\dimexpr\linewidth-6.4mm\relax}{\RaggedRight #1}\par}

% Celda marcable dentro de la rubrica (tabla de autoevaluacion). Misma
% sangria francesa, pero pensada para vivir dentro de una columna X.
\newcommand{\milcopcion}[1]{%
\makebox[5.2mm][l]{\raisebox{-.55mm}{\textcolor{milcNegro}{\milcmarca}}}%
\parbox[t]{\dimexpr\linewidth-5.2mm\relax}{\RaggedRight #1}}

% ── Recursos visuales de la guía ──────────────────────────────────────
% Declarados en el bloque `recursos:` del YAML y emitidos por el builder.
% \guiaFigura   → imagen rasterizada o PDF (foto, carta celeste, montaje).
% \guiaDiagrama → fuente TikZ/circuitikz (.tex) incrustada como vector:
%                 es la vía para circuitos y esquemáticos editables.
% #1 = ruta relativa al .tex · #2 = pie (puede ir vacío)
\newcommand{\guiaFigura}[2]{%
\begin{center}
\includegraphics[width=0.88\linewidth,height=0.38\textheight,keepaspectratio]{#1}\\[1.5mm]
{\footnotesize\itshape\textcolor{milcNegro!75}{#2}}
\end{center}
}

\newcommand{\guiaDiagrama}[2]{%
\begin{center}
\input{#1}\\[1.5mm]
{\footnotesize\itshape\textcolor{milcNegro!75}{#2}}
\end{center}
}

\begin{document}
\thispagestyle{empty}

% ── Portada ───────────────────────────────────────────────────────────
% Antes: un rectangulo de color a pagina completa. Se veia bien en pantalla,
% pero en la fotocopiadora del colegio se comia el toner y salia gris sucio.
% Ahora la banda ocupa el tercio superior y el resto va en blanco.
\begin{tikzpicture}[remember picture,overlay]
  \path[fill=milcGradePrimary]
    (current page.north west) rectangle ([yshift=-10.6cm]current page.north east);

  \node[anchor=north west,text=milcGradeText] at ([xshift=2.0cm,yshift=-1.55cm]current page.north west)
    {\sffamily\bfseries\fontsize{12}{14}\selectfont GRADO};
  \node[anchor=north west,text=milcGradeText,opacity=.85] at ([xshift=2.0cm,yshift=-2.35cm]current page.north west)
    {\sffamily\bfseries\fontsize{8.5}{10}\selectfont SERIE GUÍAS · TECNOLOGÍA E INFORMÁTICA};
  \node[anchor=north east,text=milcGradeText,opacity=.85,align=right] at ([xshift=-2.0cm,yshift=-1.55cm]current page.north east)
    {\sffamily\bfseries\fontsize{9}{12}\selectfont I.E. Sor María Juliana\\Cartago, Valle del Cauca};

  \node[anchor=west,text=milcGradeText] (gradonum) at ([xshift=1.85cm,yshift=-7.1cm]current page.north west)
    {\sffamily\bfseries\fontsize{112}{112}\selectfont 6};
  % El circulito de grado (°) solo aplica a las guias de grado («11°»); en
  % Bebras y Semillero el numero grande es un momento/modulo, no un grado.
  % Se posiciona relativo al nodo `gradonum`, no en coordenadas absolutas.
  \draw[milcGradeText,line width=1.6pt]
    ([xshift=2.0mm,yshift=-4.5mm]gradonum.north east) circle (.15cm);

  \node[anchor=west,text=milcGradeText,text width=11.4cm,align=left] at ([xshift=1.5cm,yshift=0.5cm]gradonum.east)
    {
     {\sffamily\bfseries\fontsize{9}{11}\selectfont\MakeUppercase{Período 1 · Comunicación e identidad digital}}\\[3mm]
     {\sffamily\bfseries\fontsize{21}{25}\selectfont\setlength{\baselineskip}{27pt}Riesgos\\[4pt]digitales ---\\[4pt]reconocer\\[4pt]la trampa\\[4pt]antes de caer}\\[3.5mm]
     {\sffamily\bfseries\fontsize{15}{18}\selectfont Guía 8-6-TIC}
    };
\end{tikzpicture}

\vspace*{9.4cm}
\vfill

{\sffamily\bfseries\fontsize{8}{10}\selectfont\textcolor{milcGradeDark}{LO QUE TE LLEVAS HOY}}

\begin{tcolorbox}[enhanced,colback=white,colframe=milcNegro,arc=2mm,boxrule=1.6pt,
  left=5mm,right=5mm,top=4mm,bottom=4mm,before skip=3pt,after skip=12pt,
  fontupper=\RaggedRight\fontsize{11}{15.5}\selectfont]
Una \textbf{tabla de los 4 riesgos más comunes} para chicos de tu edad, con \textbf{una señal de alerta y una acción concreta} para cada uno. Más \textbf{tu plan personal de protección} en el cuaderno: \textbf{3 adultos de confianza} con su nombre y la forma de contactarlos, \textbf{3 acciones} que harás si algo te asusta, y \textbf{una frase de cierre} firmada: \emph{``No estoy solo. Si algo en internet me asusta, le cuento a un adulto.''}.
\end{tcolorbox}

\begin{tcolorbox}[enhanced,colback=milcGris,colframe=milcGris,arc=2mm,boxrule=0pt,
  left=5mm,right=5mm,top=3.5mm,bottom=3.5mm,after skip=12pt]
{\sffamily\bfseries\fontsize{8}{10}\selectfont\textcolor{milcNegro!55}{ESTA GUÍA ES DE}}\\[3mm]
\begin{tabularx}{\linewidth}{X p{3.2cm} p{3.2cm}}
\linead{\linewidth} & \linead{3.0cm} & \linead{3.0cm}\\[-1mm]
{\fontsize{8.5}{10}\selectfont\textcolor{milcNegro!55}{Nombre completo}} &
{\fontsize{8.5}{10}\selectfont\textcolor{milcNegro!55}{Grupo}} &
{\fontsize{8.5}{10}\selectfont\textcolor{milcNegro!55}{Fecha}}\\
\end{tabularx}
\end{tcolorbox}

\vfill

\noindent\textcolor{milcLinea}{\rule{\linewidth}{1pt}}\\[2mm]
{\sffamily\bfseries\fontsize{9.5}{12}\selectfont Autor: Dr. Álvaro Cárdenas Orozco}\\[1mm]
{\sffamily\fontsize{8.5}{11}\selectfont\textcolor{milcNegro!55}{Metodología MILC · apertura ancestral · pensamiento computacional · triángulo Dussel–estoicismo–Floridi}}

\newpage

\section*{Apertura: Riesgos digitales --- reconocer la trampa antes de caer}

\begin{softbox}{milcGradeDark}{milcGradeSoft}{Saber ancestral}
\textbf{Tu abuela decía: ``en la calle, mejor caminas con alguien''.} Cuando tu mamá era niña, antes de salir del barrio le decían cosas que parecen del campo pero protegen: \emph{``no le abras a desconocidos''}, \emph{``si alguien te ofrece dulce, no aceptes''}, \emph{``si te pierdes, busca un policía o una señora''}, \emph{``avísale a alguien antes de irte''}. No eran reglas para asustarla; eran \textbf{señales para reconocer una trampa antes de caer}. Las personas que han vivido más, conocen las trampas y nos las enseñan. En internet pasa igual. \textbf{Hay trampas digitales} --- la mayoría se pueden ver venir si sabes las señales. La diferencia con la calle es que en internet la trampa puede llegar sin que tú salgas de tu cuarto. Por eso esta clase es de las más importantes del periodo.
\end{softbox}

\begin{softbox}{milcTurquesa}{milcGris}{Saber contemporáneo}
Los \textbf{riesgos digitales} son situaciones de internet que pueden hacerte daño. \textbf{No son raras:} le pasan a chicos de tu edad cada semana en Cartago y en toda Colombia. Hay 4 que vas a aprender a reconocer hoy: \textbf{ciberacoso} (cuando otros chicos se ensañan contigo en línea), \textbf{grooming} (cuando un adulto se hace pasar por amigo para hacerte daño), \textbf{retos peligrosos} (challenges que pueden lastimarte), y \textbf{estafas} (engaños para robarte dinero o datos). \textbf{Lo importante no es asustarte}, es que sepas: (1) las 5 señales que se repiten en todos los riesgos, (2) qué hacer cuando aparecen, y (3) que \textbf{nunca tienes que enfrentarlo solo}. Siempre hay un adulto de confianza que sabe ayudarte.
\end{softbox}

\begin{softbox}{milcMostaza}{milcGris}{Pregunta-puente}
\textit{Si un compañero que no conoces te empezara a escribir todos los días por TikTok preguntándote cosas personales, ¿\textbf{cómo sabrías} si quiere ser tu amigo o si está intentando algo malo? ¿Hay señales que se pueden ver?}
\end{softbox}

\begin{softbox}{milcVerde}{milcGris}{Saber y saber hacer}
Al terminar la clase: (1) podrás \textbf{identificar} los 4 riesgos digitales más comunes; (2) sabrás \textbf{explicar} las 5 señales universales de alerta; (3) habrás \textbf{aplicado} las señales a 4 escenas reales; (4) habrás \textbf{creado} tu plan personal de protección con 3 adultos de confianza.
\end{softbox}



\small
\begin{tabularx}{\textwidth}{>{\bfseries}p{3.0cm}Y}
\rowcolor{milcGradePrimary!12}Autor & Dr. Álvaro Cárdenas Orozco\\
DBA & Reconoce principios y conceptos propios de la tecnología, relacionando artefactos con su utilización segura (MEN, T\&I 6°)\\
Referentes & Saber ancestral del pregonero · Pensamiento computacional inicial · Cuidado de la palabra\\
Producto final & Una \textbf{tabla de los 4 riesgos más comunes} para chicos de tu edad, con \textbf{una señal de alerta y una acción concreta} para cada uno. Más \textbf{tu plan personal de protección} en el cuaderno: \textbf{3 adultos de confianza} con su nombre y la forma de contactarlos, \textbf{3 acciones} que harás si algo te asusta, y \textbf{una frase de cierre} firmada: \emph{``No estoy solo. Si algo en internet me asusta, le cuento a un adulto.''}.\\
\end{tabularx}
\normalsize

\puente{Hoy haces 4 cosas: identificas los 4 riesgos, aprendes las 5 señales, las aplicas a escenas reales, y firmas tu plan.}

\milcestacion{milcGradeDark}{Ruta MILC: así vamos a aprender}
\begin{tabularx}{\textwidth}{>{\centering\arraybackslash}X>{\centering\arraybackslash}X>{\centering\arraybackslash}X>{\centering\arraybackslash}X}
\phasecard{1}{milcVerde}{milcGris}{Escucha\\[-1mm]\small Observo, escucho y pregunto.} &
\phasecard{2}{milcTurquesa}{milcGris}{Entender\\[-1mm]\small Sistematizo y ordeno.} &
\phasecard{3}{milcMagenta}{milcGris}{Hacer\\[-1mm]\small Construyo y aplico.} &
\phasecard{4}{milcVino}{milcGris}{Cierre\\[-1mm]\small Evalúo y reflexiono.}
\end{tabularx}


\puente{Empezamos leyendo 4 escenas que pueden suceder en cualquier celular.}

\milcestacion{milcVerde}{Estación 1 de 3 · Escucha}

\begin{softbox}{milcVerde}{milcGris}{Escena para escuchar}
\textbf{Actividad 1 · ANALIZA --- 4 escenas reales de chicos de tu edad} (12 min · individual). Lee con calma las 4 escenas. Para cada una, en el cuaderno escribe: (a) \textbf{qué riesgo es} (ciberacoso, grooming, reto, estafa), y (b) \textbf{qué haría yo}. \textbf{Escena A:} Sara, 12 años, recibe en TikTok mensajes de un grupo de su salón burlándose de cómo se peinó. Le mandan memes feos cada día. \textbf{Escena B:} Daniel, 11 años, jugando Roblox, se hace amigo de \emph{``Lucas''}, otro jugador. Lucas dice tener 13 años y le pide foto en uniforme del colegio para \emph{``conocerlo mejor''}. \textbf{Escena C:} En un grupo de WhatsApp, le pasan a Lucía un video de un reto: aguantar la respiración hasta desmayarse. Le dicen que si no lo hace, es \emph{``floja''}. \textbf{Escena D:} A Carlos le llega un correo que dice ser de \emph{``Roblox oficial''}: \emph{``Felicitaciones, ganaste 10 mil Robux. Ingresa tu contraseña aquí para reclamar.''}.
\end{softbox}

{\sffamily\bfseries\fontsize{8}{10}\selectfont\textcolor{milcNegro!55}{MARCA CUANDO LO HAYAS HECHO}}
\milccheckrow{Leí las 4 escenas con calma.}
\milccheckrow{Identifiqué qué tipo de riesgo era cada una.}
\milccheckrow{Escribí qué haría yo, sin esperar consejo del compañero.}

\begin{infoband}{milcVerde}


\textbf{Lo que va al cuaderno (Actividad 1 · ANALIZA).} Título: «Actividad 1 --- 4 escenas y sus riesgos». \textbf{Formato:} 4 bloques (uno por escena), cada uno con 2 líneas: tipo de riesgo + qué harías. \textbf{Extensión:} 8 líneas en total. \textbf{Sabes que terminaste cuando} podrías compararle tu análisis a un compañero y defenderlo con razones.
\end{infoband}

\puente{Después de las escenas, aprendes las 5 señales que se repiten en cualquier trampa.}

\milcestacion{milcTurquesa}{Estación 2 de 3 · Entender}

\begin{softbox}{milcTurquesa}{milcGris}{Lo que vas a entender}
Las 4 escenas que leíste son distintas, pero \textbf{comparten 5 señales} que se repiten en casi cualquier riesgo digital. Si aprendes a reconocer las 5, vas a detectar trampas que ni tú sabías que existían.
\end{softbox}

\milcpaso{1}{milcTurquesa}{\textbf{Señal 1 --- Te pide secreto.} \emph{``No le digas a nadie''}, \emph{``Esto queda entre nosotros''}, \emph{``Tus papás no entenderían''}. Toda persona que te pide secreto sobre lo que pasa en internet está intentando algo. La gente honesta no necesita secretos. \textbf{Señal 2 --- Te pide algo que cruza un límite.} Foto de tu cuerpo, tu contraseña, encontrarse en persona, mucho dinero, datos privados. Cualquier persona que pida cosas que sientes raras, está cruzando un límite.}
\milcpaso{2}{milcTurquesa}{\textbf{Señal 3 --- Mueve tus emociones a propósito.} Te halaga mucho (\emph{``eres especial''}), te asusta (\emph{``si no lo haces te pasa X''}), te culpa (\emph{``por tu culpa estoy triste''}), te aísla (\emph{``tus amigos no te entienden''}). Manejarte por las emociones es la trampa. \textbf{Señal 4 --- Es muy rápido.} Una persona normal se conoce despacio. Si en una semana ya te dice \emph{``somos almas gemelas''} o te insiste mucho en cosas serias, eso es señal. La gente honesta da tiempo.}
\milcpaso{3}{milcTurquesa}{\textbf{Señal 5 --- Te aleja de los adultos de confianza.} \emph{``Si le cuentas a tu mamá, me boqueas''}, \emph{``Los profes son aburridos''}, \emph{``No le digas, te vas a meter en problemas''}. Quien quiere protegerte NUNCA pide que te alejes de tu familia. \textbf{Las 3 acciones de oro cuando aparece una señal.} (a) \textbf{Bloquea} a la persona o sal de la situación. (b) \textbf{Guarda evidencia}: captura de pantalla, mensajes guardados. (c) \textbf{Cuéntale a un adulto} de tu lista de confianza.}
\milcpaso{4}{milcTurquesa}{\textbf{4 errores que cometen los chicos de tu edad cuando aparece una trampa.} (1) Pensar \emph{``Yo no soy tonto, yo no caigo''} --- las trampas están diseñadas para gente lista. (2) Esperar a que se solucione solo --- nunca pasa. (3) Culparse \emph{``yo tuve la culpa''} --- nunca. La culpa es de quien hace daño. (4) No contar por miedo a quitar el celular --- ningún adulto serio te quita el celular por contar; te quita el celular si NO cuentas.}

\begin{drawbox}{milcTurquesa}{Las 5 señales de alerta universales}
\textbf{1.} Te pide secreto. \\[1mm]
\textbf{2.} Te pide algo que cruza un límite. \\[1mm]
\textbf{3.} Mueve tus emociones a propósito. \\[1mm]
\textbf{4.} Va muy rápido. \\[1mm]
\textbf{5.} Te aleja de los adultos de confianza. \\[1mm]
\textbf{Acciones de oro:} bloquea + guarda evidencia + cuéntale a un adulto.
\end{drawbox}

\begin{softbox}{milcMostaza}{milcGris}{Errores comunes y su corrección}
\textbf{Mitos sobre los riesgos digitales que escucharás.} (1) \emph{``Eso solo le pasa a las niñas''} --- Falso. A los niños les pasa igual o más. (2) \emph{``Solo le pasa a chicos que ya tienen problemas''} --- Falso. Cualquier chico bueno puede ser víctima. (3) \emph{``Si cuento, me van a quitar el celular''} --- Falso. Los adultos sensatos te dan más confianza si cuentas, no menos. (4) \emph{``Los retos son juegos''} --- Falso. Hay retos que han llevado a la cuenta el hospital y peor. (5) \emph{``Yo lo manejo solo''} --- No. Estos asuntos los manejan adultos con experiencia. Tú das la voz; ellos hacen el resto.
\end{softbox}

\begin{infoband}{milcTurquesa}


\textbf{Lo que va al cuaderno (Actividad 2 · EVALÚA).} Título: «Actividad 2 --- Qué tan peligrosa es cada escena y qué haría». \textbf{Formato:} 2 partes. Parte A: las 5 señales numeradas con \emph{tu propia explicación de qué significa cada una} (1 línea por señal). Parte B: las 3 acciones de oro escritas con tus palabras. \textbf{Veredicto final (3 líneas):} de las 4 escenas leídas (Sara, Daniel, Lucía, Carlos), ¿cuál es la \textbf{más peligrosa} y por qué? Justifica con 2 razones usando las 5 señales aprendidas. \textbf{Extensión:} 8 líneas + veredicto. \textbf{Sabes que terminaste cuando} tu veredicto distingue gravedad real de incomodidad pasajera.
\end{infoband}

\puente{Con las señales claras, escribes tu plan de protección personal.}

\milcestacion{milcMagenta}{Estación 3 de 3 · Hacer}

\begin{softbox}{milcMagenta}{milcGris}{Cómo vas a construir tu producto}
\textbf{Actividad 3 · CREA --- Tu plan personal de protección} (25 min · individual). Vas a hacer en una hoja del cuaderno tu \textbf{plan personal}: los \textbf{3 adultos de confianza} a los que les contarías si algo pasa, sus datos para contactarlos, y \textbf{3 acciones específicas} que harías si una señal de alerta aparece. Este plan se queda en tu cuaderno, es solo tuyo, y te lo aprendes para que cuando lo necesites no tengas que pensar.
\end{softbox}

\milcpaso{1}{milcMagenta}{\textbf{Paso 1 --- Mis 3 adultos de confianza.} En la parte superior de la hoja escribe el título: \textbf{``Mis 3 adultos de confianza''}. Después haz una tabla de 3 filas y 3 columnas. Columna 1: nombre. Columna 2: relación contigo (mamá, papá, profe, abuela, tía, coordinadora). Columna 3: cómo lo contacto (celular, donde lo veo cada día, su correo). \textbf{Llena las 3 filas} con personas reales.}
\milcpaso{2}{milcMagenta}{\textbf{Paso 2 --- Reglas para escoger.} Los 3 adultos deben cumplir esto: (a) que tú \emph{de verdad} confíes en él o ella (no solo que sea adulto), (b) que sepa escucharte sin gritarte de entrada, (c) que tenga forma rápida de localizarse. Si alguno no cumple, cámbialo por otro.}
\milcpaso{3}{milcMagenta}{\textbf{Paso 3 --- Mis 3 acciones cuando aparece una señal.} Debajo de la tabla escribe el subtítulo: \textbf{``Mis 3 acciones''}. Lista numerada con: (1) \emph{``Lo primero: respiro, no respondo, no envío fotos''}; (2) \emph{``Capturo pantalla y guardo el chat sin borrar''}; (3) \emph{``Le cuento a uno de mis 3 adultos hoy mismo, hoy mismo''}. Puedes escribirlas en tus palabras pero las 3 ideas tienen que estar.}
\milcpaso{4}{milcMagenta}{\textbf{Paso 4 --- Mis líneas de ayuda.} Después de tus 3 acciones, escribe \textbf{3 líneas oficiales gratuitas} que existen en Colombia para denunciar: \emph{``Te Protejo Colombia: teprotejo.org''}, \emph{``Policía CAI Virtual: 123''}, \emph{``ICBF: 141 (Bienestar Familiar)''}. Las anotas aunque nunca las uses --- saber que existen ya protege.}
\milcpaso{5}{milcMagenta}{\textbf{Paso 5 --- Frase de cierre.} En el centro inferior de la hoja, escribe en grande la frase: \emph{``No estoy solo. Si algo en internet me asusta, le cuento a un adulto.''}. Subráyala. Esta frase es como tu mantra cuando pase algo difícil.}
\milcpaso{6}{milcMagenta}{\textbf{Paso 6 --- Firma.} En la esquina inferior derecha, firma con tu nombre completo y la fecha de hoy. Este plan se queda en el cuaderno. La idea es que NO lo tengas que volver a leer hasta que lo necesites --- pero cuando lo necesites, esté ahí.}

\begin{drawbox}{milcMagenta}{Antes de cerrar tu plan}
\checkbox La tabla tiene 3 adultos con sus 3 datos cada uno. \\[1mm]
\checkbox Los 3 adultos cumplen las reglas (confianza, escucha, contacto rápido). \\[1mm]
\checkbox Están las 3 acciones de oro escritas. \\[1mm]
\checkbox Están las 3 líneas de ayuda anotadas (Te Protejo, CAI, ICBF). \\[1mm]
\checkbox La frase de cierre está subrayada. \\[1mm]
\checkbox Está firmado con nombre y fecha.
\end{drawbox}

\begin{softbox}{milcMagenta}{milcGris}{Plantilla de trabajo}
\textbf{Estructura del plan.} \textit{Arriba:} Título. \textit{Centro arriba:} Tabla de 3 adultos. \textit{Centro:} Mis 3 acciones. \textit{Centro abajo:} 3 líneas oficiales. \textit{Centro inferior:} frase subrayada. \textit{Esquina:} firma.
\end{softbox}

\begin{infoband}{milcMagenta}


\textbf{Lo que va al cuaderno (Actividad 3 · CREA).} Título: «Actividad 3 --- Mi plan personal de protección». \textbf{Formato:} hoja completa del cuaderno con tabla de adultos, lista de acciones, líneas oficiales y frase firmada. \textbf{Extensión:} 1 hoja entera. \textbf{Sabes que terminaste cuando} sabrías exactamente qué hacer y a quién acudir si en este momento te pasara algo en una app.
\end{infoband}

\puente{El producto son las 4 escenas analizadas + el plan personal con 3 adultos.}

\milcestacion{milcMostaza}{Producto de la sesión}
\textbf{Plan personal de protección · 3 adultos · 3 acciones · firmado}.
\begin{softbox}{milcMostaza}{milcGris}{Criterios de calidad}
(1) Los 3 adultos son personas reales con datos correctos de contacto. (2) Las 3 acciones están claras y son ejecutables. (3) Están las 3 líneas oficiales de ayuda con su forma de contacto. (4) La frase de cierre está visible y subrayada. (5) Está firmado y fechado.
\end{softbox}

\puente{Antes de salir, te aseguras de tener los teléfonos y formas de contactar a los 3 adultos.}

\milcestacion{milcVino}{Evaluación liberadora · cinco dimensiones}

\begin{softbox}{milcVino}{milcGris}{Sentido de la evaluación}
Evaluar no es solo revisar una nota. Es reconocer qué aprendí, qué cambió en mí, qué emociones manejé, cómo me relaciono con la ciudad y la comunidad, y qué le dejaré a quien viene detrás.
\end{softbox}

\textbf{Mi reflexión liberadora en cinco dimensiones.} Completa cada frase iniciada:

\small
\begin{tabularx}{\textwidth}{>{\bfseries\RaggedRight\arraybackslash}p{3.4cm}Y>{\RaggedRight\arraybackslash}p{4.6cm}}
\rowcolor{milcVino}\color{white}Dimensión & \color{white}\bfseries Mi respuesta (frame starter) & \color{white}\bfseries Algo que descubrí\\
1. Personal & Lo que más cambió en mí fue \linead{6.5cm} & \linead{4.2cm}\\
2. Emocional & La emoción más fuerte fue \linead{2.5cm} y la regulé \linead{3cm} & \linead{4.2cm}\\
3. Ciudadana & En democracia esto importa porque \linead{6cm} & \linead{4.2cm}\\
4. Local (Cartago) & En mi barrio sería útil para \linead{6.5cm} & \linead{4.2cm}\\
5. Intergeneracional & A mi abuela/abuelo le diría \linead{6.5cm} & \linead{4.2cm}\\
\end{tabularx}
\normalsize

\begin{infoband}{milcVino}
\textbf{Para la próxima sustentación.} Vuelve a esta tabla en seis meses. Verás que las respuestas cambian — no porque lo aprendido cambie, sino porque tú cambias. La evaluación liberadora es una conversación contigo a través del tiempo.
\end{infoband}


\puente{Cerramos con 3 voces que te ayudan a entender por qué pedir ayuda es de valiente, no de débil.}

\milcestacion{milcGradeDark}{Triángulo de pensamiento · cierre filosófico}

\begin{softbox}{milcVino}{milcGris}{Dussel · lente del nosotros}
\textit{No nacimos para enfrentar todo solos. Ser persona es saber pedir ayuda y darla.}

Hay una mentira común: que pedir ayuda es de cobardes y aguantarse solo es de valientes. \textbf{Es al revés.} Los adultos que más respetas en tu vida --- tu papá, tu abuela, tu profe --- pidieron ayuda cuando fue necesario. Si no lo hubieran hecho, no estarían bien hoy. \textbf{Pedir ayuda es lo más valiente que puedes hacer cuando algo te asusta.} Especialmente en internet, donde a veces lo que parece pequeño esconde algo grande.

\textbf{¿Estoy creyendo el mito de que ``si pido ayuda parezco débil''? ¿O ya entendí que pedir ayuda es lo más valiente?}
\end{softbox}
\linead{\linewidth}\\[1mm]
\linead{\linewidth}

\begin{softbox}{milcMostaza}{milcGris}{Estoicismo · Séneca · cuidado interior}
\textit{Lo que te asusta de noche, en silencio, pierde fuerza cuando se lo cuentas a otra persona.}

Cuando algo te asusta en el celular y te lo guardas, \textbf{crece}. Cuando se lo cuentas a un adulto, \textbf{se reduce} a su tamaño real. Casi siempre lo que parecía gigante se vuelve manejable. La diferencia entre crecer y reducirse depende de un solo paso: contarlo.

\textbf{¿Estoy cargando algo solo en silencio? Si lo cuento, ¿quién está listo para escucharme?}
\end{softbox}
\linead{\linewidth}\\[1mm]
\linead{\linewidth}

\begin{softbox}{milcTurquesa}{milcGris}{Floridi · lente de la infoesfera}
\textit{Tu mejor protección en internet no es saber más tecnología que el adversario. Es tener una red humana detrás de ti.}

Algunos chicos creen que la mejor protección es \emph{``saber más que los malos''} de tecnología. Es bueno, pero no es suficiente. \textbf{Tu mejor protección es tu red humana}: mamá, papá, profe, hermana, abuela. La gente que te conoce a ti. Los hackers más expertos del mundo siguen teniendo redes humanas detrás. Tú también deberías.

\textbf{¿Estoy construyendo mi red humana, o creo que puedo manejar todo solo con mi celular?}
\end{softbox}
\linead{\linewidth}\\[1mm]
\linead{\linewidth}

\begin{infoband}{milcNegro}
\textbf{Nota para el docente.} Las tres formulaciones del triángulo son la síntesis del autor de la guía, no citas textuales de los pensadores. Si vas a citarlos en clase, remite a la obra original.
\end{infoband}

\milcestacion{milcVino}{Mi autoevaluación · marca una casilla por fila}

{\small
\setlength{\extrarowheight}{2.2pt}
\begin{tabularx}{\textwidth}{>{\RaggedRight\arraybackslash\bfseries}p{3.0cm}YYY}
\rowcolor{milcVino}\color{white}\bfseries Criterio & \color{white}\bfseries Lo logré & \color{white}\bfseries En proceso & \color{white}\bfseries Necesito apoyo\\[.6mm]
Comprensión del tema & \milcopcion{Explico las ideas con mis palabras.} & \milcopcion{Reconozco las ideas, falta precisión.} & \milcopcion{Necesito repasar lo básico.}\\[.8mm]
\rowcolor{milcGris}Pensamiento computacional & \milcopcion{Usé los pilares al armar mi producto.} & \milcopcion{Usé algunos pilares.} & \milcopcion{Necesito releer la estación 2.}\\[.8mm]
Aplicación y producto & \milcopcion{Mi producto cumple los criterios.} & \milcopcion{Está armado, con ajustes pendientes.} & \milcopcion{Aún está en construcción.}\\[.8mm]
\rowcolor{milcGris}Reflexión 5 dimensiones & \milcopcion{Respondí cada una con honestidad.} & \milcopcion{Respondí algunas con detalle.} & \milcopcion{Mis respuestas son muy generales.}\\[.8mm]
Triángulo de pensamiento & \milcopcion{Conecté las 3 voces con mi trabajo.} & \milcopcion{Conecté al menos una.} & \milcopcion{No las relaciono con mi proceso.}\\[.8mm]
\end{tabularx}}

\vspace{3mm}
\textbf{Hoy entendí que:}\\[1mm]
\linead{\linewidth}\\[1mm]
\linead{\linewidth}

\textbf{Mi compromiso para la próxima sesión es:}\\[1mm]
\linead{\linewidth}\\[1mm]
\linead{\linewidth}

\begin{softbox}{milcMostaza}{milcGris}{Cita de despedida · Marco Aurelio}
\textit{``El obstáculo en el camino se vuelve el camino. Lo que estorba al camino se vuelve camino.''}\\[1mm]
Cada error de hoy, bien observado, se convierte en pieza del camino que pisarás mañana.
\end{softbox}

\small
\begin{tabularx}{\textwidth}{>{\bfseries}p{3.6cm}Y}
\rowcolor{milcGradeDark!12}Nombre completo: & \linead{10cm}\\
Fecha: & \linead{4cm} \quad Curso/grupo: \linead{4cm}\\
Mi firma: & \linead{9cm}\\
\end{tabularx}
\normalsize

\begin{infoband}{milcGradeDark}
\textbf{Para la próxima sesión.} Esta semana habla con UNO de tus 3 adultos de confianza. Dile: \emph{``profe (mamá/abuela/tía), te quería avisar que si alguna vez en internet me pasa algo raro, te voy a contar a ti.''}. No tiene que pasar nada, solo dejar claro que existe el camino. La próxima clase aprenderás sobre \textbf{comunicación responsable}: cómo escribir y hablar en línea de manera que sumes en lugar de hacer daño. Hoy aprendiste a recibir; la próxima clase aprendes a emitir bien.
\end{infoband}

\end{document}
